% Minimal retained preamble for lane-p2 cards taken from
% sources/logic-fl02-conservation-of-ramsey-s-theorem-f/pinned-original-source.tex
% (Conservation of Ramsey's theorem for pairs, arXiv:2402.11616v3).
%
% Copied verbatim from the author's own file:
%   line 2      : \documentclass{article}
%   lines 16-25 : the author's theorem declarations
%   lines 46-88 : the author's macro definitions for the reverse-mathematics principles
%   line 91     : \card
%   lines 104-107: the natural sums and products, which rely on mathabx (source line 103)
%   line 116    : \qt
%
% Omitted author preamble lines and why (nothing was renamed or redefined):
%   line 3 (graphicx), 7 (url), 9 (cancel): no figure, url or cancel use in the retained spans;
%   line 6 (cleveref): the assembler adds its own cleveref machinery, a second load clashes;
%   line 8 (\usepackage[utf8x]{inputenc}): inputenc with the utf8x option fights the forced
%       xelatex engine;
%   lines 28-37 (the rep@theorem machinery and \newreptheorem{maintheorem}): the retained spans
%       use no rep-theorem environment;
%   lines 40-43 (xcolor and the three author comment macros): not used by the retained spans;
%   lines 93-99 (the LS2 font encoding block for \smashtimes based on stix): the host TeX Live has
%       no stix directory and the author himself never uses the result -- line 105 maps \ntimes to
%       \dottimes, and line 110 shows \smashtimes already commented out in the source.
% mathabx is retained because \dotplus, \dottimes and \dotdiv come from it.
\documentclass{article}
\usepackage{amsmath,amssymb,amsthm}
\usepackage{mathabx}
\usepackage{cancel}
\usepackage{cleveref}

\newtheorem{theorem}{Theorem}
\numberwithin{theorem}{section}
\newtheorem{maintheorem}[theorem]{Main Theorem}
\newtheorem{lemma}[theorem]{Lemma}
\newtheorem{question}[theorem]{Question}
\newtheorem{proposition}[theorem]{Proposition}
\newtheorem{statement}[theorem]{Statement}
\newtheorem{remark}[theorem]{Remark}
\newtheorem{definition}[theorem]{Definition}
\newtheorem{corollary}[theorem]{Corollary}

\newcommand{\RCA}{\mathsf{RCA}}
\newcommand{\WKL}{\mathsf{WKL}}
\newcommand{\SRT}{\mathsf{SRT}}
\newcommand{\EM}{\mathsf{EM}}
\newcommand{\SEM}{\mathsf{SEM}}
\newcommand{\ADS}{\mathsf{ADS}}
\newcommand{\SCAC}{\mathsf{SCAC}}
\newcommand{\SADS}{\mathsf{SADS}}
\newcommand{\COH}{\mathsf{COH}}
\newcommand{\ACA}{\mathsf{ACA}}
\newcommand{\qvdash}{\operatorname{{?}{\vdash}}}
\newcommand{\nqvdash}{\operatorname{{?}{\nvdash}}}
\newcommand{\BSig}{\mathsf{B}\Sigma^0}
\newcommand{\ISig}{\mathsf{I}\Sigma^0}
\newcommand{\LSig}{\mathsf{L}\Sigma^0}
\newcommand{\BPi}{\mathsf{B}\Pi^0}
\newcommand{\IPi}{\mathsf{I}\Pi^0}
\newcommand{\LPi}{\mathsf{L}\Pi^0}
\newcommand{\IDel}{\mathsf{I}\Delta^0}
\newcommand{\BCDe}{\mathsf{BC}\Delta^0}
\newcommand{\WF}{\mathsf{WF}}
\newcommand{\RT}{\mathsf{RT}}
\newcommand{\BME}{\mathsf{BME}}
\newcommand{\TT}{\mathsf{TT}}
\renewcommand{\O}{\mathcal{O}}
\renewcommand{\L}{\mathcal{L}}
\newcommand{\C}{\mathcal{C}}
\newcommand{\R}{\mathcal{R}}
\renewcommand{\S}{\mathcal{S}}
\newcommand{\T}{\mathcal{T}}
\renewcommand{\P}{\mathcal{P}}
\newcommand{\M}{\mathcal{M}}
\newcommand{\Nc}{\mathcal{N}}
\renewcommand{\P}{\mathcal{P}}
\newcommand{\Q}{\mathcal{Q}}
\newcommand{\NN}{\mathbb{N}}
\newcommand{\PP}{\mathbb{P}}
\newcommand{\uh}[0]{{\upharpoonright}}
\newcommand{\card}{\operatorname{card}}
\newcommand{\nplus}{\dotplus}
\newcommand{\ntimes}{\dottimes}
\newcommand{\nminus}{\dotdiv}
\newcommand{\nsum}{\dot{\sum}}
\def\qt#1{``#1''}
\title{Ramsey's theorem for pairs preserves a prescribed ordinal well-foundedness while computing the jump}
\author{Quentin Le Houérou, Ludovic Levy Patey, Keita Yokoyama}
\date{arxiv:2402.11616v3; CC BY 4.0}
\begin{document}
\maketitle
\noindent\textit{Editorial note.} The counted unit is the article's own Proposition prop:rt22-wf at source lines 1356-1363 together with the proof environment that begins on the next line (1364-1370); the pairing is the ordinary adjacent one.\par\medskip
\noindent\textit{Editorial note.} The three context spans are the statements of the COH, SEM and SADS propositions that the counted proof applies in sequence; their proofs are not retained and they are counted by sibling cards of this lane. Note that the SADS proposition at source lines 1102-1109 is retained here only as context: as a counted unit that interval is dispositioned, because its own proof environment contains the counted lemma lem:sads-forcing-question-sigma01.\par\medskip
\noindent\textit{Editorial note.} The counted proof carries the author citation of Fiori-Carones et al inside its first line. The author's own bibliography file biblio.bib from the pinned e-print of arXiv:2402.11616v3 is attached to this task unedited and complete, with the bibtex backend, so that the citation resolves.\par\medskip
\noindent\textit{Editorial note.} The retained preamble reproduces verbatim the author's document class at source line 2, his theorem declarations at source lines 16-25, his macro definitions at source lines 46-88 and 91, and his natural-operation macros at source lines 104-107, and it retains mathabx and cleveref; the stix-dependent LS2 block and the inputenc load with the utf8x option are not retained.\par\medskip
\noindent\textit{Editorial note.} The article's conservation theorems and its other propositions, lemmas and corollaries are not counted here.\par\medskip
\noindent\textit{Allowed standard prerequisites (editorial).}\begin{itemize}
\item Omega-extension of a topped model into a model of WKL\_0-star. The proof starts from a result of Fiori-Carones et al (their Lemma 3.2) that from the auxiliary set P one obtains a P-computable countable coded model of WKL\_0-star containing the jump of Y, hence a set P\_0 with P Turing above P\_0 Turing above the jump of Y; the citation is carried in the proof and its entry is attached with this card.. It supplies the intermediate set P\_0 that the first application of the cohesive-set proposition needs.
\item The three forcing propositions chained by the proof. The article proves that the COH forcing produces a cohesive set while preserving well-foundedness of a power of alpha and controlling P\_0 relative to the jump; that the SEM forcing over that model produces a transitive subtournament with the same properties for one lower power; and that the SADS forcing produces a homogeneous set preserving well-foundedness of alpha itself; their statements are retained with this card at source lines 461-468, 760-767 and 1102-1109.. The proof applies the three propositions one after another to descend from well-foundedness of omega\_4 to the fourth power of alpha down to alpha itself, and to obtain the homogeneity of the final set.
\item Stable colourings, SEM instances and SADS instances. A colouring on the model is an instance of the cohesiveness principle, and its restriction to a cohesive set is stable; a stable colouring restricted further is an instance of the ascending-descending sequence principle on a transitive set; the article identifies these instances in the proof.. These identifications are what allow the three propositions to be applied in sequence to the same colouring.
\end{itemize}\medskip
\begin{proposition}\label[proposition]{prop:coh-wf}
Consider a countable model $\M = (M,S) \models \RCA_0 + \BSig_2 + \WF(\alpha)$ topped by a set $Y \in S$, where $\alpha \leq \epsilon_0$ is an ordinal. Then for every countable sequence of sets $\vec{R} \in S$ and every~$P \gg Y'$ such that $\M[P] \models \RCA_0^*$, there exists $G \subseteq M$ such that 
\begin{enumerate}
    \item $G$ is $\vec{R}$-cohesive ;
    \item $P \geq_T (G \oplus Y)'$ ;
    \item $\M[G] \models \RCA_0 + \BSig_2 + \WF(\alpha)$.
\end{enumerate}
\end{proposition}

\begin{proposition}\label[proposition]{prop:sem-wf}
Consider a countable model $\M = (M,S) \models \RCA_0 + \BSig_2 + \WF(\sup_{\beta < \alpha} \omega^{\beta \times \omega})$ topped by a set $Y \in S$, where~$\alpha \in M$ is a non-zero ordinal $\leq \epsilon_0$. Then for every infinite stable tournament $R$ in $S$ and every set~$P \gg Y'$ such that $\M[P] \models \RCA_0^*$, there exists $G \subseteq M$ such that  
\begin{enumerate}
    \item $G$ is an $M$-infinite transitive subtournament ;
    \item $P \geq_T (G \oplus Y)'$ ;
    \item $\M[G] \models \RCA_0 + \BSig_2 + \WF(\alpha)$.
\end{enumerate}
\end{proposition}

\begin{proposition}\label[proposition]{prop:sads-wf}
Consider a countable model $\M = (M,S) \models \RCA_0 + \BSig_2 + \WF(\sup_{\beta < \alpha} \omega^{\beta \times \omega})$ topped by a set $Y \in S$, where $\alpha \in M$ is an ordinal $\leq \epsilon_0$. Then for every linear order $L = (M, \leq_L) \in S$ of order type $\omega + \omega^*$, there exists $G \subseteq M$ such that  
\begin{enumerate}
    \item $G$ is an $M$-infinite ascending or descending sequence for $\leq_L$;
    \item $Y' \geq_T (G \oplus Y)'$ ;
    \item $\M[G] \models \RCA_0 + \BSig_2 + \WF(\alpha)$.
\end{enumerate}
\end{proposition}

\begin{proposition}\label[proposition]{prop:rt22-wf}
Consider a countable model $\M = (M,S) \models \RCA_0 + \BSig_2 + \WF(\omega_4^{\alpha})$ topped by a set~$Y \in S$, where~$\alpha \in M$ is an ordinal~$\leq \epsilon_0$. Then, for every coloring~$f : [M]^2 \to 2$ in~$S$ and every set~$P \gg Y'$ such that $\M[P] \models \RCA_0^*$, there exists $G \subseteq M$ such that 
\begin{enumerate}
    \item $G$ is an $M$-infinite $f$-homogeneous set ;
    \item $P \geq_T (G \oplus Y)'$ ;
    \item $\M[G] \models \RCA_0 + \BSig_2 + \WF(\alpha)$.
\end{enumerate}
\end{proposition}

\begin{proof}
By Fiori-Carones et al~\cite[Lemma 3.2]{fiori2021isomorphism}, there exists a $P$-computable countable coded model of~$\WKL_0^*$ containing~$Y'$. In particular, there exists a set~$P_0 \subseteq M$ such that $P \gg P_0 \gg Y'$.
Let~$\vec{R} = R_0, R_1, \dots$ be defined by $R_x = \{ y : f(x, y) = 1 \}$.
By \Cref{prop:coh-wf} applied to~$\M$ and $P_0$, there exists some $M$-infinite $\vec{R}$-cohesive set~$G_0 \subseteq M$ such that $P_0 \geq_T (G_0 \oplus Y)'$ and $\M[G_0] \models \RCA_0 + \BSig_2 + \WF(\omega_3^{\alpha})$. In particular, $f : [G_0]^2 \to 2$ is stable. One can see~$f : [G_0]^2 \to 2$ as an instance of~$\SEM$.
By \Cref{prop:sem-wf} applied to~$\M[G_0]$ and $P$, there exists some~$M$-infinite $f$-transitive set~$G_1 \subseteq G_0$ such that $P \geq_T (G_1 \oplus G_0 \oplus Y)'$ and $\M[G_1 \oplus G_0] \models \RCA_0 + \BSig_2 + \WF(\omega_2^\alpha)$. One can see $f : [G_1]^2 \to 2$ as an instance of~$\SADS$.
By \Cref{prop:sads-wf} applied to~$\M[G_0 \oplus G_1]$ and $P$, there exists some~$M$-infinite $f$-homogeneous set~$G \subseteq G_1$ such that $P \geq_T (G \oplus G_1 \oplus G_0 \oplus Y)'$ and $\M[G \oplus G_1 \oplus G_0] \models \RCA_0 + \BSig_2 + \WF(\alpha)$.
\end{proof}

\bibliographystyle{plain}
\bibliography{retained-biblio}
\end{document}
